% !PRINTED 245
\weberpage{245}
%<<CONT
in which the $a_{h,k}$ are integers, and the conditions
\[
a_{h,k} = a_{k,h}, \quad a_{p,p} = 0
\]
are satisfied.

\begin{center}
\textbf{§ 70.\quad \gesperrt{Transformation Equations of the First Level.}}
\end{center}

The invariant equation is of great theoretical importance, partly because
of its general validity (for even $n$ as well), partly because of the ease
with which the linear transformations can be applied to $j(\omega)$. The
actual computation of this equation, however, turns out to be so complicated, and
the numerical coefficients are so large, that the computation has so far
been carried out only in the simplest case $p=2$. On the other hand, by
making use of other modular functions, one can obtain far simpler
transformation equations.

Concerning the principle to be applied here we remark the following:

If there is given any system of $\nu$ functions of $\omega$, corresponding
to the $\nu$ systems of transformation numbers $a$, $c$, $\partial$, which
we shall denote by
\weq{1}{\Phi_{a,c,\partial}}
and which are permuted among one another \gesperrt{isomorphically} with the functions
\weq{2}{j\left(\frac{c+\partial\omega}{a}\right)}
by the substitutions
\weq{3}{\left(\omega, -\frac{1}{\omega}\right), \quad (\omega, \omega+1)}
then (1) is expressible \gesperrt{rationally} through (2) and through $j(\omega)$; for
the function

\weq{4}{F_n[x,j(\omega)] \sum^{a,c,\partial} \frac{\Phi_{a,c,\partial}}{x-j\left(\frac{c+\partial\omega}{a}\right)} = \Psi[x,j(\omega)]}

remains unaltered under the substitutions (3), and is therefore (for an
indeterminate $x$) a \gesperrt{rational function} of $j(\omega)$, and moreover a
polynomial in $x$ of degree at most $\nu-1$;
$F_n[x,j(\omega)]$ has the same meaning as in (17) of the preceding
section, and $\nu=\psi(n)$ is the degree of this function with respect to
$x$. From (4), however, on setting
\[
x = j\left(\frac{c+\partial\omega}{a}\right)
\]
one obtains:
%CONT>>
